\clearpage
\section{Learned refinement pilot}\label{app:training}

The pilot trained a learned correction to the analytic decoder using the fitting-only inputs of Appendix~\ref{app:evidence}. All logged optimizer updates completed without a nonfinite skip. The corrected objective ran through the final curriculum phase. Its control ended earlier, so comparisons late in training cannot isolate the effect of the objective or the number of updates.

Training loss and its retention term fell as the corrected runs progressed, but the checkpoint transfer screens and the release decision use independent evaluations on synthetic data. The checkpoint that met the selection criteria on wide DEV was evaluated once on fresh row splits from development tables. It scored slightly below the signed release on that schedule, with an interval covering zero, and also missed the task and held-out MLP guards. The learned tier stayed a pilot. The plotted aggregates and their input and output digests are in \path{evidence/addendum-k/} of the scoreboard.

\begin{figure}[!htbp]
  \centering
  \includegraphics[width=\linewidth]{\BudgetData/training-k5r2.pdf}
  \caption{Training and downstream transfer in the learned-tier pilot. Panels A and B show window medians of training loss and its retention term. Vertical lines mark curriculum boundaries. Solid traces indicate the corrected objective and dashed traces indicate the control. Panel C shows paired Transfer Score differences against the signed release at \VThree{PrimaryMult}$\times$, with \VThree{CiLevel} family-bootstrap intervals. Each group orders runs by seed and each panel uses its own row-split schedule. The dotted line marks the I2 selection margin. I1 and I2 selected checkpoints. I3 evaluated the selected checkpoint on fresh DEV row splits. The plotted points cannot establish a training-duration effect.}
  \label{fig:training}
\end{figure}
